\section{讲义式导读：如何判断一个 AI for Math 系统是否靠谱}

学习本期时，可以把“AI for Math 系统是否靠谱”拆成六个问题。第一，它能否准确理解数学对象和定理陈述，而不是只做自然语言 paraphrase。第二，它能否把非形式证明转成 Lean/Coq 等 proof assistant 可检查的形式。第三，它能否在证明失败时读取 goal state，并修复而不是重写一段幻觉文本。第四，它能否利用定理库，而不是每次从零开始。第五，它能否给数学家带来研究速度提升，而不只是做已知题。第六，它是否能解释证明为什么有意义。

\begin{importantbox}{六问法}
一个 AI for Math 系统若只满足前两问，它更像 formalization 工具；若能满足第三、第四问，它开始具备 Agent 能力；若能满足第五、第六问，它才接近数学研究伙伴。
\end{importantbox}

\subsection{为什么这六问重要}

它们把“AI 会不会数学”从模糊问题拆成可检查能力。理解对象对应语义能力；形式化对应语言与类型系统能力；读取 goal state 对应环境反馈；利用定理库对应检索和记忆；提升研究速度对应产品价值；解释意义对应数学审美。高标准 AI for Math 需要这些能力共同出现。

\subsection{本章小结}

这六问会贯穿后文。读者不必先懂 Lean 细节，也能用它判断不同 AI for Math 路线到底是在做题、做证明，还是在尝试进入真正研究。

